\section{Binary boundary tensors for the bracket}
\label{sec:tensors}

\subsection{Orient each circle instead of recording its connectivity}

Work initially in a commutative extension of
$\Z[A,A^{-1}]$ containing an invertible element $\rho$ with
\begin{equation}\label{eq:rho-fourth}
 \rho^4=-A^2.
\end{equation}
Then $\rho^4+\rho^{-4}=\delta$. By Theorem~\ref{thm:turning}, the sum of the
weights of the two orientations of any smoothing circle is exactly
$\delta$. Thus every factor $\delta^{c(s)}$ in the bracket state sum can
be evaluated by independently orienting the circles.

Choose a canonical dart $d_e$ on each diagram edge; the implementation
chooses the smaller dart index. A binary edge state gives that edge an
arrow, either in the direction of $d_e$ or against it. Write
$\epsilon_d=+1$ when the arrow points out of the vertex at dart $d$, and
$\epsilon_d=-1$ otherwise. Thus $\epsilon_{\alpha(d)}=-\epsilon_d$.
A smoothing is compatible with these arrows if each of its pairs has one
incoming and one outgoing dart.

For a compatible smoothing $s$ at crossing $v$, let $r_v(s,\epsilon)$ be
the sum of its two local turns. It belongs to $\{-2,0,2\}$. Define the
local crossing tensor and the edge weight by
\begin{align}
 T_v(\epsilon)&=\sum_{\substack{s\in\{0,1\}\\s\text{ compatible}}}
          A^{1-2s}\rho^{r_v(s,\epsilon)},\label{eq:crossing-tensor}\\
 E_e(\epsilon)&=\rho^{\epsilon_{d_e}b_{d_e}}.\label{eq:edge-tensor}
\end{align}
The crossing tensor vanishes unless two darts point in and two point out.
Only six of the sixteen arrow configurations can occur. The two alternating
configurations each support both smoothings; the others support one.

\begin{theorem}[Exact two-state contraction]
\label{thm:binary-tensor}
For a validated nonempty classical knot diagram,
\begin{equation}\label{eq:network}
 \mathcal B_D(A)=
 \sum_{\text{one binary arrow per edge}}
     \prod_v T_v(\epsilon)\prod_e E_e(\epsilon).
\end{equation}
For any crossing order, a direct frontier contraction of this expression
has at most $2^{w_i}$ keys after a prefix whose cut has $w_i$ diagram edges.
No connectivity or disk condition on the prefix is required.
\end{theorem}

\begin{proof}
Expand every sum in \eqref{eq:crossing-tensor}. A nonzero term consists of
a smoothing state and an arrow on each edge such that the arrows extend
coherently through every smoothing arc. This is exactly an independent
choice of orientation for each smoothing circle. Its weight is the
smoothing monomial $A^{n-2|s|}$ times $\rho$ raised to the sum of all local
and edge turns. Summing the orientations of one circle gives
$\rho^4+\rho^{-4}=\delta$. Doing this independently for all circles recovers
\eqref{eq:bracket-state-sum}.

To contract a prefix, sum all arrow variables whose two incident crossings
are processed and retain those with exactly one processed endpoint. The
result is a tensor indexed by the latter binary variables. There are
$2^{w_i}$ assignments. A loop at one crossing has both endpoints processed
together and is summed locally. This is ordinary distributivity in a
finite product of local tensors and does not use a boundary matching.
\end{proof}

\subsection{A continuation statement with its exact scope}

Let $P$ be a set of processed crossings, let $Q$ be its complement, and
fix the same full-diagram turning data. Assign each edge factor to its
canonical endpoint. The contraction of $P$ produces a vector $u_P$ and
that of $Q$ produces the complementary vector $u_Q$ on the same cut
assignments. Their contraction is $\mathcal B_D$.

This gives an exact linear summary for the bracket under the remaining
contraction. It does not imply that the summary preserves the relative
Khovanov complex, nor that it distinguishes every topological tangle.
The distinction explains how the binary state bound can coexist with
large matching sets and with the baseline's lower bounds on explicit
categorified checkpoints. Decategorification can identify data that a
chain-complex computation must retain.

\subsection{One cyclotomic phase per frontier key}

For exact arithmetic we later set $A=B^2$ and $\rho=\zeta B$, where
$\zeta^4=-1$. At first sight every coefficient is a four-coordinate
integer in $\Z[\zeta]$. A stronger statement removes most of that storage.

\begin{lemma}[A vertex potential modulo four]
\label{lem:phase-potential}
For an oriented edge from slot $j$ at vertex $v$ to slot $k$ at vertex $u$,
there is a function $\nu$ on vertices with values in $\Z/4\Z$ such that
\begin{equation}\label{eq:phase-potential}
 b_d+j-k-2=\nu(u)-\nu(v)\pmod4.
\end{equation}
\end{lemma}

\begin{proof}
The expression on the left is antisymmetric modulo four. Around a face,
the next departing slot equals the preceding arriving slot plus one,
modulo four. Consequently $\sum(j-k)=|f|\pmod4$. The face sum of the
left side of \eqref{eq:phase-potential} is therefore
$\sum b_d-|f|=\pm4=0\pmod4$. Apply
Lemma~\ref{lem:closed-cochain} over $\Z/4\Z$ and take its vertex potential.
\end{proof}

\begin{theorem}[Phase determined by the boundary]
\label{thm:one-phase}
Fix a set $P$ of processed crossings and a binary assignment on its cut.
All nonzero monomials contributing to that tensor entry have the same
power of $\zeta$ modulo four. The entry therefore has the form
$c\zeta^r$, with $c\in\Z$ and $0\le r<4$, after the integer
specialization and the nonnegative shifts described below.
\end{theorem}

\begin{proof}
At a crossing, the local turning exponent satisfies
\[
 r_v=\sum_{d=(v,j)}\epsilon_d j\pmod4.
\]
Indeed, an arc from incoming $j$ to outgoing $k$ turns by
$k-j-2\pmod4$; summing two arcs subtracts four.
Every allowed crossing has two incoming and two outgoing darts, so
$\sum_{d\text{ at }v}\epsilon_d=0$. We may consequently add
$\nu(v)\sum\epsilon_d$ without changing the phase.

An internal edge of the prefix contributes, from its two endpoints and
its edge factor,
\[
 \epsilon_d\bigl(b_d+j-k+\nu(v)-\nu(u)\bigr)
       =2\epsilon_d=2\pmod4
\]
by Lemma~\ref{lem:phase-potential}. This contribution is independent of
its arrow. A cut edge contributes only terms determined by its fixed
arrow, its processed endpoint, and the fixed convention assigning its
edge factor. Thus the total phase is fixed by the boundary assignment and
by $P$. The extra powers of $B$ used for denominator clearing have no
$\zeta$ phase. Powers differing by four have opposite signs because
$\zeta^4=-1$, and those signs are absorbed into the integer $c$.
\end{proof}

The implementation stores this phase together with one signed integer.
When terms reach the same key, it checks that their phases agree before
adding their integers. The check is an implementation invariant backed by
the theorem, rather than an assumption that discards other coordinates.
Zero entries are removed. A closed result has phase zero, since every
smoothing circle contributes a multiple of four.

\subsection{The executable contraction}

The contraction uses a deterministic crossing order. At each step it
performs the following operations.
\begin{enumerate}
\item Identify the current cut edges incident to the crossing, cut edges
that survive untouched, newly opened edges, and self-loops.
\item Compile the local arrow table from
\eqref{eq:crossing-tensor}--\eqref{eq:edge-tensor}, summing self-loop
orientations immediately. At most eight smoothing--orientation
contributions occur before equal monomials are aggregated.
\item For each stored boundary key, read its incident arrow bits and
apply the compatible local transitions. Copy the untouched bits and
append the new bits to form the next key.
\item Multiply by the compiled local monomials, merge equal keys using
Theorem~\ref{thm:one-phase}, and remove zero entries.
\end{enumerate}
The local tensors have bounded arity. An arbitrary number of holes in the
processed region changes neither the arity nor the binary dimension of a
cut edge. The implementation's cap counts actual tensor updates, not the
number of smoothing states of the full diagram.

\begin{remark}[Why this does not improve Khovanov automatically]
The evaluation in \eqref{eq:network} is an identity for a Laurent
polynomial. Its signed sums can cancel information that would survive in
homology. Applying the same data reduction to a differential without
proving a homotopy or homology comparison would be invalid. A useful next
problem is to identify precisely which additional data, if any, can be
retained at controlled cost; Section~\ref{sec:research} discusses this.
\end{remark}
